\documentclass{report}
\usepackage{amsmath,amssymb}
\setlength{\parindent}{0mm}
\setlength{\parskip}{1em}
\begin{document}
\begin{center}
\rule{6in}{1pt} \
{\large Steven Hamilton \\
{\bf Coupled Radiation Transport and Thermomechanics using the AMP and Denovo Codes}}

Oak Ridge National Laboratory \\ P O Box 2008 \\ Oak Ridge \\ TN 37831-6170
\\
{\tt hamiltonsp@ornl.gov}\\
Kevin Clarno\\
Bobby Philip\\
Rahul Sampath\\
Mark Berrill\\
Srikanth Allu\\
Mark Baird\\
	Jim Banfield\end{center}

The ability to accurately model the behavior of nuclear fuels during
irradiation in a nuclear reactor remains one of the most significant
impediments to the development of improved reactor designs. Modeling the
coupled effects of heat transfer and structural mechanics on nuclear
fuels in a reactor environment is vital to the development of predictive
fuel performance capabilities. Existing methodologies utilize `reference
pin' approximations in which a single fuel pin is selected to be
representative of a much larger set of pins. This approach, however, is
unable to account for assembly-level effects due to interactions between
pins. In this study we demonstrate the use of a heat generation source
produced by solving the Boltzmann neutron transport equation on a full
light water reactor (LWR) assembly with thermomechanics calculations
performed on every pin, capturing effects that are generally ignored. The
transport equation is solved using the Denovo 3D discrete ordinates code
and the thermomechanics solutions use the AMP multiphysics package.
Results from computations performed on the Jaguar supercomputer on more
than 40,000 processors will be presented.


\end{document}
